% Unit 2 free-response set -- 2 long (10) + 3 short (4) = 32 pts.
\documentclass{apchem-exam}

\apcunit{2}
\apcunittitle{Compound Structure and Properties}
\apcdoctitle{Free-Response Set}
\apcsubtitle{CED 2.1--2.7 \textbullet\ 5 questions \textbullet\ 32 points \textbullet\ 50 minutes}

\begin{document}
\apctitleblock

\begin{apcnote}[Scope --- two CED exclusions apply]
Covers \ced{2.1} through \ced{2.7}, worth
\SI{7}{\percent}--\SI{9}{\percent} of the exam.

Topic \ced{2.3} excludes \textbf{specific crystal structures}, so
unit-cell geometry is enrichment and is not asked. Topic \ced{2.7}
excludes \textbf{hybridization involving $d$ orbitals} --- for more than
four electron domains only the \emph{shape} is required, never an
\ce{sp^3d} label --- and also excludes \textbf{molecular orbital theory}.
Neither appears below.

Every structural explanation must reach the \textbf{Coulombic} level.
Rubrics award zero for ``it wants a full octet''.
\end{apcnote}

\examsection{II}{Free Response}{5 questions \textbullet\ 32 points \textbullet\ 50 minutes}{\calculatorok}

% =====================================================================
\frq{long}{10}{\ced{2.5} \ced{2.6} \ced{2.7} \textbullet\ Lewis structure, formal charge, shape and polarity}

Consider the thiocyanate ion, \ce{SCN-}, whose skeleton is
\ce{S-C-N}.

\begin{parts}
  \item Determine the total number of valence electrons and draw a Lewis
        structure. \workspace[3cm]
        \keyonly{\begin{apcanswer}Valence electrons:
        $6(\ce{S}) + 4(\ce{C}) + 5(\ce{N}) + 1(\text{charge}) =
        \mathbf{16}$.

        The best structure is \ce{[S=C=N]-}: two double bonds, with two
        lone pairs on sulfur and two on nitrogen. Carbon is central and
        carries no lone pair.\end{apcanswer}}
  \item Assign the formal charge on each atom in your structure.
        \workspace[2.8cm]
        \keyonly{\begin{apcanswer}Formal charge $=$ valence electrons
        $-$ lone-pair electrons $-\tfrac{1}{2}$(bonding electrons).

        \ce{S}: $6 - 4 - \tfrac{1}{2}(4) = \mathbf{0}$ \quad
        \ce{C}: $4 - 0 - \tfrac{1}{2}(8) = \mathbf{0}$ \quad
        \ce{N}: $5 - 4 - \tfrac{1}{2}(4) = \mathbf{-1}$

        The charges sum to $-1$, matching the ion's overall
        charge.\end{apcanswer}}
  \item An alternative structure places a triple bond between C and N.
        State which structure is preferred and justify using formal
        charge. \answerlines[3]{Both contribute, but the preferred
        structure is the one placing the negative formal charge on the
        \emph{more electronegative} atom. In \ce{[S=C=N]-} the $-1$ sits
        on nitrogen; the triple-bonded form \ce{[S-C#N]-} puts $-1$ on
        sulfur, which is less electronegative. The arrangement that leaves
        negative charge on the more electronegative atom is favoured.}
  \item Predict the shape and bond angle about the central carbon, and
        state the model used. \answerlines[2]{Two electron domains on
        carbon and no lone pair, so by VSEPR the ion is \textbf{linear}
        with a bond angle of \SI{180}{\degree}.}
\end{parts}

\begin{rubric}
  \rpoint{3}{(a) 1 pt 16 valence electrons \emph{including} the one for
    the charge; 1 pt a structure with a complete octet on every atom;
    1 pt all lone pairs shown with brackets and charge}
  \rpoint{3}{(b) 1 pt the formal-charge relationship stated or used
    correctly; 1 pt all three values; 1 pt noting they sum to the ion
    charge}
  \rpoint{2}{(c) 1 pt identifying the more electronegative atom as the one
    that should carry the negative charge; 1 pt applying it to choose
    between the structures}
  \rpoint{2}{(d) 1 pt linear; 1 pt \SI{180}{\degree} from two domains with
    no lone pair}
\end{rubric}

% =====================================================================
\frq{long}{10}{\ced{2.1} \ced{2.2} \textbullet\ the bonding continuum and potential energy}

\begin{parts}
  \item Bonding is better described as a continuum than as three separate
        categories. Place \ce{NaCl}, \ce{HCl} and \ce{Cl2} on that
        continuum and justify using electronegativity.
        \workspace[3cm]
        \keyonly{\begin{apcanswer}\ce{Cl2}: two identical atoms, so the
        electronegativity difference is \textbf{zero} and the pair is
        shared equally --- a \textbf{pure covalent} bond.

        \ce{HCl}: a moderate difference, so the pair is shared but pulled
        toward chlorine --- \textbf{polar covalent}, with partial charges
        $\delta+$ on H and $\delta-$ on Cl.

        \ce{NaCl}: a large difference, large enough that the electron is
        effectively transferred rather than shared --- \textbf{ionic}.

        The three differ in \emph{degree} of electron-pair sharing, not in
        kind.\end{apcanswer}}
  \item Sketch, in words, how the potential energy of two approaching
        atoms varies with internuclear distance, and identify what the
        minimum represents. \workspace[3.2cm]
        \keyonly{\begin{apcanswer}At large separation the energy is
        essentially zero --- the atoms barely interact. As they approach,
        attraction between each nucleus and the other's electrons lowers
        the energy, and the curve falls.

        At very short distance the two \emph{nuclei} repel one another
        strongly and the curve turns sharply upward.

        The \textbf{minimum} is the balance point of those two effects.
        Its horizontal position is the \textbf{bond length}; its depth
        below zero is the \textbf{bond energy}, the energy released on
        forming the bond and required to break it.\end{apcanswer}}
  \item The \ce{Cl2} bond energy is \SI{243}{\kilo\joule\per\mole} and
        that of \ce{Br2} is \SI{193}{\kilo\joule\per\mole}. Explain the
        difference. \answerlines[3]{Bromine atoms are larger, so the
        bonded nuclei sit further apart and the shared pair is further
        from both nuclei. The Coulombic attraction holding the bond
        together is correspondingly weaker, giving a longer, weaker bond.}
  \item State how the potential-energy curve for \ce{Br2} differs from
        that for \ce{Cl2}. \answerlines[2]{Its minimum is
        \textbf{shallower} (smaller bond energy) and lies at a
        \textbf{greater} internuclear distance (longer bond).}
\end{parts}

\begin{rubric}
  \rpoint{3}{(a) 1 pt each substance correctly placed with an
    electronegativity-difference justification}
  \rpoint{3}{(b) 1 pt attraction lowering the energy on approach; 1 pt
    nuclear--nuclear repulsion at short range; 1 pt the minimum read as
    bond length and bond energy}
  \rpoint{2}{(c) 1 pt larger atoms give a longer bond; 1 pt weaker
    Coulombic attraction over that greater distance. ``\ce{Br2} is less
    stable'' earns 0}
  \rpoint{2}{(d) 1 pt shallower minimum; 1 pt at greater distance}
\end{rubric}

% =====================================================================
\frq{short}{4}{\ced{2.3} \textbullet\ ionic solids}

\begin{parts}
  \item Explain why ionic solids are hard and brittle, while being poor
        conductors as solids but good conductors when molten.
        \workspace[3.2cm]
        \keyonly{\begin{apcanswer}\textbf{Hard:} every ion is held by
        strong electrostatic attraction to all its oppositely charged
        neighbours throughout the lattice.

        \textbf{Brittle:} a blow displaces one plane of ions by one
        position, bringing ions of \emph{like} charge face to face. The
        attraction along that plane is replaced by strong repulsion and
        the crystal splits.

        \textbf{Conductivity:} conduction needs mobile charges. In the
        solid the ions are locked in fixed positions; melting frees them
        to move, so the liquid conducts.\end{apcanswer}}
\end{parts}

\begin{rubric}
  \rpoint{4}{1 pt hardness from strong electrostatic attraction; 1 pt
    brittleness from like charges being brought adjacent; 1 pt ions fixed
    in the solid; 1 pt ions mobile in the melt. Note \ced{2.3} excludes
    specific crystal structures --- do not require a unit-cell description}
\end{rubric}

% =====================================================================
\frq{short}{4}{\ced{2.4} \textbullet\ metals and alloys}

\begin{parts}
  \item Explain, using the electron sea model, why metals conduct
        electricity and can be drawn into wire.
        \answerlines[3]{A metal is a lattice of cations in a sea of
        \textbf{delocalized} electrons belonging to no particular atom.
        Those electrons drift through the solid under a potential
        difference, so the metal conducts. When the metal is drawn out,
        layers of cations slide over one another while the electron sea
        flows with them and keeps binding the lattice together, so it
        deforms instead of fracturing.}
  \item Steel is an alloy of iron with a small amount of carbon and is
        much harder than pure iron. Give the structural reason.
        \answerlines[2]{The carbon atoms are a different size and occupy
        positions between the iron atoms, disrupting the regular layers.
        Layers can no longer slide over one another easily, so the metal
        resists deformation and is harder.}
\end{parts}

\begin{rubric}
  \rpoint{2}{(a) 1 pt delocalized electrons free to move; 1 pt layers
    slide while the electron sea maintains bonding}
  \rpoint{2}{(b) 1 pt differently sized atoms disrupt the lattice; 1 pt
    layers slide less easily}
\end{rubric}

% =====================================================================
\frq{short}{4}{\ced{2.7} \textbullet\ hybridization and bond counting}

\begin{parts}
  \item State the hybridization of each carbon in ethanoic acid,
        \ce{CH3COOH}. \answerlines[2]{The \ce{CH3} carbon has four
        electron domains: \ce{sp^3}. The \ce{COOH} carbon has three
        domains (one \ce{C=O}, one \ce{C-O}, one \ce{C-C}): \ce{sp^2}.}
  \item Count the sigma and pi bonds in the whole molecule.
        \answerlines[3]{\textbf{Seven sigma:} three \ce{C-H}, one
        \ce{C-C}, one \ce{C=O} (its first component), one \ce{C-O} and one
        \ce{O-H}. \textbf{One pi:} the second component of the \ce{C=O}.
        A double bond is always one sigma plus one pi.}
\end{parts}

\begin{rubric}
  \rpoint{2}{(a) 1 pt each carbon, requiring the domain count behind it}
  \rpoint{2}{(b) 1 pt seven sigma; 1 pt one pi. Note \ced{2.7} excludes
    $d$-orbital hybridization --- never require an \ce{sp^3d} label}
\end{rubric}

\examtotal

\end{document}
